\documentclass[10pt,a4paper]{article}
\usepackage[utf8]{inputenc}\usepackage[T1]{fontenc}\usepackage[english]{babel}
\usepackage{graphicx,float,booktabs,amsmath,amssymb,siunitx,xcolor,tabularx,enumitem}
\usepackage[margin=2.0cm]{geometry}\usepackage[hidelinks]{hyperref}
\usepackage{newunicodechar}\newunicodechar{σ}{$\sigma$}\newunicodechar{–}{--}\newunicodechar{−}{$-$}\newunicodechar{≈}{$\approx$}
\usepackage{textcomp}
\graphicspath{{bilder/}{../../docs/bilder/}}
\hypersetup{pdfauthor={Thomas Schmiegelt},pdftitle={Strength review package: coil support structure of a stellarator}}
\newcommand{\bild}[3]{\begin{figure}[H]\centering\includegraphics[width=#2\linewidth]{#1}\caption{#3}\end{figure}}
\newcommand{\lesehilfe}[1]{\begin{quote}\small\textit{How to read the figure:} #1\end{quote}}
\newcommand{\kasten}[2]{\par\smallskip\noindent\fcolorbox{gray}{gray!8}{\parbox{0.96\linewidth}{\textbf{#1} #2}}\par\smallskip}
\title{\textbf{Strength Review Package}\\[2pt]\large Coil support structure of a stellarator -- model, results and
eight questions for structural analysts}
\author{Thomas Schmiegelt \textbar{} hobby project \texttt{stellarator-design}}
\date{3 November 2026}
\begin{document}\maketitle

\kasten{Hobby project, non-commercial.}{No connection to any customer project, research institute or company, not
peer-reviewed. Calculations, code and draft text were produced with the help of AI agents; the human set goal and
direction. All numbers come from simplified models (linear elastic, rigid bearings, magnetic load only). That is
exactly why I am asking for your critical view: the questions are in Section~\ref{sec:fragen}.}

\noindent\textbf{Audience.} Structural analysts (FEM, codes and standards) without a fusion background. Everything
fusion-specific is briefly explained. Reading time about 20 minutes; for the questions alone,
Sections~\ref{sec:kurz} and \ref{sec:fragen} are enough.

\section{At a glance}\label{sec:kurz}
\begin{itemize}[itemsep=1pt]
\item \textbf{Component:} casings of steel 1.4429 (316LN) around 24 superconducting coils, joined by steel
connectors, supported on an inner and outer support ring. Load: Lorentz force during magnet operation, up to
\SI{1.1}{GN} per coil.
\item \textbf{Criterion:} σ99 = 99\% value of the von Mises stress over the steel volume; limit \SI{500}{MPa}
($\tfrac23 R_{p0.2}$ at \SI{4}{K}, ITER grade C3).
\item \textbf{Result:} of 199 evaluated designs in a study, 7 meet all limits. Recommended: design 0169 with
σ99 = \SI{308}{MPa}, also on a finer mesh (\SI{319}{MPa}).
\item \textbf{Main issue:} the peak values sit at the connector junctions and grow with every mesh refinement. For
four of five candidates σ99 rises by 34--58\% on a finer mesh.
\item \textbf{New (3 Nov):} in 0169 the connectors were so close that they almost formed a plate. Re-computed:
separate webs are mesh-stable (+1\%), a continuous plate is not (+47\%).
\end{itemize}

\section{What it is about}
A stellarator confines a plasma hotter than 100 million degrees in a twisted magnetic field. The field comes from
superconducting coils with an irregular, three-dimensional shape. The design here: \SI{1}{GW} net electric power,
major radius $R_0=\SI{10.4}{m}$, field $B_0=\SI{6.5}{T}$, high-temperature superconductor (REBCO), operation at about
\SIrange{4}{20}{K}. Each coil carries about \SI{20}{MA}. In its own field and that of its neighbours it feels the
Lorentz force $f=J\times B$: \SIrange{50}{140}{MN/m} along the coil, in total up to \SI{1.1}{GN} per coil towards
the machine centre. Wanted: a structure of coil casings, connectors and bearings that carries these forces without
straining the superconductor by more than \SI{0.3}{\%}.

\bild{modell3d_0169_schnitt_en}{0.72}{Design 0169 in a stepped cut (3D model in PicoGK).}
\lesehilfe{From inside out: plasma (pink), first wall and blanket (yellow), shield (green), vacuum vessel (grey),
then the 24 steel coil casings (blue) with the superconducting winding pack (orange). Between the casings sit the
connectors (dark blue). The inner support cylinder and the outer ring are only drawn: they show where the bearing
forces go, but are neither computed nor designed.}

\section{Model}
\begin{description}[itemsep=1pt,style=nextline]
\item[Geometry] Coil centre lines from the optimisation (simsopt \cite{simsopt}). One hexahedral mesh per coil along
its path: 64 stations for 24 coils; in the cross-section 2 wall cells, \SI{2}{cm} potting and $4\times4$ winding-pack
cells. Connectors as hexahedral blocks between the casings.
\item[Materials] Casing 1.4429 (X2CrNiMoN17-13-3, AISI 316LN), $E=\SI{200}{GPa}$. Winding pack homogenised. Potting
quartz-sand epoxy, $E=\SI{10}{GPa}$. Connectors steel, $E$ as a design variable \SIrange{20}{200}{GPa}.
\item[Load] Lorentz force density from all coils (Biot--Savart), magnet operation only.
\item[Bearings] Patches on the casing outer surface at the smallest and largest radius, \textbf{rigidly clamped}.
\item[Solver] Elmer \cite{elmer}, linear elastic, linear 8-node hexahedra (Elmer type 808), MUMPS.
\item[Chain] Plasma $\to$ coils $\to$ package $\to$ strength $\to$ temperature, 1--3 minutes per design. A study of
200 designs varied coil count, coil--plasma distance, connector count and stiffness, wall thicknesses and bearing
radius.
\end{description}

\section{How the variants come about: computed automatically, rejected automatically}\label{sec:auto}
No design was drawn by hand. A chain of programs turns eleven design variables into coils, package, strength model,
temperatures and a 3D model, and assesses the result against eight limits. A Latin hypercube plan spreads the
variables evenly over their ranges. The AI wrote and started the chain; the 200 designs ran unattended in 8.5 hours.

\bild{ablauf_en}{0.98}{The automatic chain of the design study and what remains of it.}
\lesehilfe{Top: the stations every design passes through. Bottom left, the funnel: of 200 planned designs one
failed in the solver (aborted without result), in 11 the casings intersect (not buildable). Of the remaining 188, 7 meet all limits and
3 more miss one narrowly. Bottom right: why the 181 rejected designs fail, i.e.\ the limit with the smallest margin.
Mostly it is the steel stress, but almost half fail on something else (field accuracy, superconductor strain,
neutron fluence). A pure strength design would not filter them out.}

\bild{galerie_en}{0.98}{16 rejected designs, evenly spread over the ranking (3D models in PicoGK).}
\lesehilfe{Top left the design furthest from the limits, bottom right the one that fails most narrowly. Above each
image: number, coil count, σ99 and the limit it fails. "Margin" is the distance to the limit relative to the limit
($-100$\% means twice the allowable). Red frame: clearly missed, yellow: narrowly. High stress alone does not decide:
0021 has only \SI{558}{MPa}, but its superconductor would receive a multiple of the allowable neutron fluence. The
catalogue in the variant report shows all 188 buildable designs like this.}

\bild{film_feld}{0.82}{Still from the film of all variants (second 30): the rejected designs sink.}
\lesehilfe{Each ring is a complete design from the study. Red: clearly missed, yellow: narrowly, blue: admissible.
The inadmissible ones sink in turn, worst first; at the end the seven blue ones remain. The film is in the repository
(\texttt{docs/filme/varianten.mp4}).}

\section{How the structure came about}
The structure was not prescribed but derived from a force analysis. The first model reached \SI{2129}{MPa}. Instead
of varying parameters, we first examined how the force arises.
\begin{enumerate}[itemsep=1pt]
\item \textbf{Force along the coil:} \SIrange{50}{140}{MN/m}, net per coil up to \SI{1.1}{GN} inwards.
\item \textbf{Rope or beam:} a curved segment carries transverse load as tension ($T=f_N/\kappa$); a straight one is
bent. The largest load (\SI{\approx130}{MN/m}) sits on the almost straight inner leg ($\kappa<\SI{0.08}{1/m}$ over
about \SI{7}{m}).
\item \textbf{Hand calculation:} simply supported beam, $q\approx\SI{130}{MN/m}$, $l\approx\SI{7}{m}$:
$M=ql^2/8\approx\SI{800}{MNm}$; hollow section $H=\SI{1.36}{m}$, $h=\SI{0.755}{m}$, $W\approx\SI{0.38}{m^3}$, hence
$\sigma_B\approx\SI{2100}{MPa}$. FEM gave \SI{2129}{MPa}: the stress came from bending of the inner leg.
\item \textbf{Flaw in the initial model:} supported only on the outside, while the force pulls inwards.
\item \textbf{Precedents:} in a tokamak the straight leg rests against a bucking cylinder (compression instead of
bending); in ARIES-CS a continuous shell passes the transverse load from coil to coil.
\end{enumerate}

\bild{festigkeit_tragwerk_en}{0.95}{Same 16 coils, five support structures A--E.}
\lesehilfe{Top: four key figures per structure; bottom: the stress seen from above (blue low, red high). A: bearings
outside, 6 connectors per pair. B: bearings inside ($-32$\%, $-30$ to $-50$\% was expected). C: inside and outside.
D: 12 stiff connectors per pair (shorter free bending length, $M\propto l^2$). E: additionally thickest walls. The
stress drops from \SI{2129}{} to \SI{481}{MPa}, steel grows from \SI{15.6}{} to \SI{37}{kt}. Since then
"inside and outside" is standard; connectors, walls and bearing radius are design variables of the study.}

\section{Criteria and limits}\label{sec:grenzen}
\begin{table}[H]\centering\small\begin{tabular}{@{}lll@{}}\toprule
Quantity & Limit & Origin \\\midrule
σ99 in casing steel & \SI{500}{MPa} & $S_m=\tfrac23R_{p0.2}$ at \SI{4}{K}, ITER grade C3 ($R_{p0.2}\ge\SI{750}{MPa}$) \cite{mitchell}\\
Membrane stress $P_m$ per wall section & \SI{500}{MPa} & linearisation through the wall thickness \\
$P_m+P_b$ & \SI{750}{MPa} & $1.5\,S_m$ \\
Winding-pack strain & \SI{0.3}{\%} & REBCO irreversible from \SIrange{0.45}{0.72}{\%} \\\bottomrule
\end{tabular}\end{table}

\bild{stahl_kennwerte_en}{0.78}{Strength of 1.4429 / 316LN at room temperature and at \SI{4}{K}.}
\lesehilfe{Blue: yield strength, green: tensile strength, yellow: resulting allowable stress. At room temperature
(left, table value \cite{schweizerfn}) only about \SI{200}{MPa} would be allowed. Measured at \SI{4.2}{K}, the yield
strength is more than three times higher (\SI{1070}{MPa} \cite{fedorikova2023}; with high nitrogen in thick plates
above \SI{1000}{MPa} \cite{sakurai2024}). The dashed line is the limit used. The casing is cold in operation; no
current flows during cool-down.}

\section{Results}
\bild{doe1_uebersicht_en}{0.85}{Design study: 199 designs, σ99 against six design variables.}
\lesehilfe{Each circle is a design, the colour gives the coil count. Filled circles meet all limits. The dashed line
is the \SI{500}{MPa} limit, the y axis is logarithmic. Only designs with 20 or 24 coils get below it. The clearest
trend is the coil--plasma distance (bottom centre): coils further away are more strongly shaped.}

\paragraph{Mesh study.} The candidates were recomputed with 1.5 times as many stations along the coil and with 3
instead of 2 wall cells.
\bild{doe1_netz_en}{0.95}{Mesh study: same coils and parameters, different mesh.}
\lesehilfe{(b) shows σ99, (c) the peak value. A finer resolution through the wall ("cross-section fine") changes
hardly anything. More stations along the coil ("1.5× stations") raise σ99 by 34--58\% for four of five candidates,
for 0169 by only 4\%. The peak value rises for all; it depends on the mesh.}

\paragraph{Location of the peaks.} 77--100\% of the 200 most highly stressed elements lie less than \SI{1}{m} from
a connector junction. For 0169 the peak rises from \SI{717}{} to \SI{1262}{MPa} on the finer mesh, σ99 only from
\SI{308}{} to \SI{319}{MPa}.

\paragraph{Categories.} The stress was linearised across each wall section into membrane ($P_m$) and membrane plus
bending ($P_m+P_b$).
\bild{doe1_kategorien_en}{0.95}{Stress categories of the admissible designs.}
\lesehilfe{(a) Peak values: they sit at the single peak at the connector end (local membrane stress $P_L$). (b) Value
met by 99\% of the wall sections: for the best designs $P_m$ stays below \SI{400}{MPa}. E (test) is the earlier
support structure test E for comparison.}

\paragraph{Connectors as webs or as a plate.} In 0169 (13 connectors per pair) neighbouring connectors were exactly
one connector width apart. They touched or overlapped, and steel was counted twice in the overlaps. Hypothesis
before the computation: the almost continuous plate makes 0169 mesh-stable. Three variants with the same coils were
computed.
\bild{verbinder_modi_en}{0.72}{Design 0169: three kinds of connectors, standard and fine mesh.}
\lesehilfe{Blue: standard mesh, orange: 1.5 times as many stations. The number above each pair is the increase due
to the finer mesh. Separate webs (spacing at least one width) are higher (\SI{387}{MPa}) but stay stable (+1\%). A
continuous plate with shared nodes jumps by 47\% to almost \SI{500}{MPa}. The hypothesis is refuted: the mesh
stability of 0169 does not come from the plate. The peaks are in well-shaped elements (scaled Jacobian 0.92--0.96),
86--96\% of them in the casing. With separate webs 0169 still remains admissible.}

\paragraph{Softer connectors.} If all connectors of 0169 are made uniformly softer, σ99 rises (+16\% at
\SI{50}{GPa}). The connectors are the relief path.

\section{Known weaknesses}
\begin{enumerate}[itemsep=1pt]
\item Bearings rigidly clamped; the support cylinder is missing as a component (optimistic).
\item Connector junctions sharp-edged, hence mesh-dependent peaks. The assessment applies to the walls, not to the
junctions.
\item Primary load only: no cool-down from \SI{293}{} to \SI{4}{K}, no dead weight, no fatigue, no fracture mechanics.
\item Winding pack homogenised, potting perfectly bonded (no debonding).
\item σ99 is not a code quantity but a figure of merit for the design search.
\item Linear hexahedra with only 2 cells through the wall thickness.
\end{enumerate}

\section{Questions for you}\label{sec:fragen}
\begin{enumerate}[itemsep=3pt]
\item \textbf{Criterion:} Is σ99 defensible as a search criterion, or should we assess linearised categories ($P_m$,
$P_L$, $P_b$) along selected sections from the start?
\item \textbf{Peaks at the junction:} The peak values at the connector end grow with the mesh (0169: \SI{717}{}
$\to$ \SI{1262}{MPa}). How do you usually handle this: hot-spot structural stress, notch stress with a fillet
radius, submodel?
\item \textbf{Junction shape:} Which shape would you choose for a steel connector between two casings (gap about
\SIrange{0.2}{1}{m}, forces in the MN range per connector): fillet, tapered thickness, longer contact face,
continuous shell?
\item \textbf{Plate or webs:} In the model the continuous plate is mesh-sensitive (+47\%), separate webs are not
(+1\%). Is this a known effect (stiffness jump at the plate edge)? Does anything speak against a continuous shell
(assembly, thermal contraction during cool-down, inspectability of welds)?
\item \textbf{Load cases:} Which are missing, in which order? Proposal: cool-down, dead weight, fast discharge
(quench), cyclic load for fatigue.
\item \textbf{Support:} How would you model the support cylinder, as a component with contact or as a spring boundary
condition? What stiffness would be realistic?
\item \textbf{Material at 4~K:} Is $\tfrac23R_{p0.2}$ (ITER magnet criteria) sufficient as a limit, or does fracture
mechanics govern at \SI{4}{K} ($K_{IC}$, crack growth)?
\item \textbf{Mesh:} How much do linear 8-node hexahedra with 2 cells through the wall underestimate bending stress
(shear locking)? Are quadratic elements (20 nodes) worth it, or is a third cell enough? Our test with 3 wall cells
hardly changed σ99.
\end{enumerate}

\section*{Documents}
Project report (strength chapter), variant report with data sheets of all admissible designs and a catalogue of all
188 buildable designs, source code (\texttt{stellarator/kopplung.py}, \texttt{hexmodell.py},
\texttt{doe\_auswertung.py}) and a film of all variants: \url{https://github.com/ThomasSchmiegelt/stellarator-design}
(MIT licence).

\begin{thebibliography}{9}\small
\bibitem{simsopt} M.~Landreman et al.: SIMSOPT: A flexible framework for stellarator optimization. J.\ Open Source
  Softw.\ 6 (2021) 3525. doi:10.21105/joss.03525.
\bibitem{elmer} Elmer FEM, CSC -- IT Center for Science, \url{https://www.elmerfem.org}.
\bibitem{mitchell} N.~Mitchell et al.: ITER TF structure, grades C1--C3 (IAEA FEC contribution IT/1-4); year and
  wording not verified.
\bibitem{schweizerfn} Strength values of steel, after Roloff/Matek, Maschinenelemente; schweizer-fn.de,
  \url{https://www.schweizer-fn.de/festigkeit/festigkeitswerte/stahl/stahl_start.php} (accessed 3 Nov 2026).
\bibitem{fedorikova2023} A.~Fedoriková, P.~Petroušek, T.~Kvačkaj, R.~Kočiško, M.~Zemko: Development of Mechanical
  Properties of Stainless Steel 316LN-IG after Cryo-Plastic Deformation. Materials 16 (2023) 6473.
  doi:10.3390/ma16196473.
\bibitem{sakurai2024} T.~Sakurai, O.~Umezawa, Y.~Ono: Strength and Fracture Toughness of Type 304 and 316 Austenitic
  Stainless Steels at 4.2~K. IOP Conf.\ Ser.: Mater.\ Sci.\ Eng.\ 1302 (2024) 012002.
  doi:10.1088/1757-899X/1302/1/012002.
\end{thebibliography}
\end{document}
